\documentclass[12pt]{article}
\usepackage[margin=1in]{geometry}
\usepackage{setspace}
\usepackage{mathptmx}
\usepackage{amsmath}
\usepackage{float}
\usepackage[section]{placeins}
\usepackage{caption}
\DeclareCaptionFont{ninept}{\fontsize{9pt}{11pt}\selectfont}
\usepackage{afterpage}
\usepackage[T1]{fontenc}
\usepackage[utf8]{inputenc}
\usepackage{textcomp}
\usepackage{changepage}
\usepackage{color}
\usepackage{bbm}
\usepackage{graphicx}
\usepackage[dvipsnames]{xcolor}
\usepackage[colorlinks,citecolor=blue,linkcolor=black,urlcolor=black]{hyperref}
\usepackage{natbib}
\usepackage{appendix}
\usepackage{chngcntr}
\usepackage{booktabs}
\usepackage{subcaption}
\usepackage{pdflscape}
\usepackage[flushleft]{threeparttable}
\usepackage{rotating}
\newcommand{\ra}[1]{\renewcommand{\arraystretch}{#1}}
\usepackage{titlesec}
% Applied Economics heading style: level 1 bold, level 2 bold italic, level 3 italic.
\titleformat{\section}{\normalfont\large\bfseries}{\thesection}{1em}{}
\titleformat{\subsection}{\normalfont\normalsize\bfseries\itshape}{\thesubsection}{1em}{}
\titleformat{\subsubsection}{\normalfont\normalsize\itshape}{\thesubsubsection}{1em}{}
\doublespacing
\begin{document}
\title{When Demand Crosses Borders: Market Entry and the Cannabis Border Premium}
\author{Mahyar Rezazadeh\thanks{PhD Candidate, Department of Economics, Simon Fraser University. Email: \href{mailto:smr20@sfu.ca}{smr20@sfu.ca}}\\ {\small \emph{Simon Fraser University}}}
\date{\today}
\maketitle
\thispagestyle{empty}
\begin{abstract}
\noindent Recreational marijuana legalization has been adopted unevenly across U.S. states, creating adjacent legal and nonlegal markets and opportunities for cross-border demand. This paper examines whether demand from neighboring nonlegal states generates a border premium in legal cannabis markets, measured by dispensary revenues. Using dispensary locations and revenue data from New Mexico from 2022-Q2 to 2025-Q4, it tests for a border premium under a relatively free-entry regime. Border ZIP codes generate about 89 percent more dispensary revenue per resident and have 46 percent greater dispensary density than interior ZIP codes, a pattern consistent with free entry. At the individual-dispensary level, however, there is no statistically significant difference in revenue between border and interior locations. Stronger demand near the border therefore attracts additional dispensaries and expands market size, while competition limits revenue gains for individual firms. Comparing New Mexico with Colorado, where local licensing restrictions constrain entry, this paper finds suggestive evidence that entry barriers attenuate the border premium, though the difference is only marginally significant. Overall, cross-border demand expands legal cannabis markets, but the extent to which markets capture it depends on the regulatory structure governing entry.
\end{abstract}

\medskip
\noindent \textbf{JEL classification:} L13, L51, H71, R12, K23 \\
\noindent \textbf{Keywords:} Cannabis legalization; Cross-border demand; Market entry regulation; Free entry; Border premium
\newpage
\section{Introduction}

Recreational marijuana is legal in a growing number of U.S. states but remains prohibited in many others. By the end of 2025, 23 states had operating recreational cannabis markets, while recreational retail sales remained prohibited in the other 27 states.

This creates a source of captive demand for border markets in legal states that are close to populations in nonlegal states. Residents of areas in nonlegal states near a legal market may be able to reach a licensed dispensary with a short drive. Those living farther from legal-state borders or dispensaries face higher travel costs and have less access to legal recreational cannabis.

There is evidence that this type of captive demand may arise when neighboring jurisdictions tax, regulate, or prohibit the same good differently. Studies of alcohol, cigarettes, gasoline, lotteries, and cannabis show that consumers travel across borders in response to differences in taxes, prices, and legal access. These responses generally weaken with distance because travel becomes more costly (\cite{AsplundFribergWilander2007, Lovenheim2008, BeattyLarsenSommervoll2009, ChandraHeadTappata2014}).

For the marijuana market, understanding whether cross-border demand translates into higher dispensary revenue is important for state and local policy decisions on licensing, zoning, and taxation. When cross-border spillovers are present, regulation may reduce consumption of a sin good and its associated social costs, but it means giving up tax revenue and local economic activity generated by consumers from neighboring nonlegal states. \cite{HansenMillerWeber2020} show that cross-border demand can disappear quickly. When Oregon legalized recreational cannabis, Washington retailers near the border lost sales because consumers no longer needed to cross state lines to buy legally. But those retailers had already entered the market before Oregon legalized. This raises a different and important question: would as many firms have entered those border markets if they had expected the extra demand to be temporary? The answer matters for entry regulation because it speaks to whether licensing restrictions should be tightened or relaxed in markets where cross-border demand may be temporary.

The setting of this paper makes it possible to study how cross-border demand affects market outcomes when firms are relatively free to enter. In \cite{HansenMillerWeber2020}, retailers in Washington were already operating when Oregon legalized recreational marijuana.
 In contrast, New Mexico implemented recreational marijuana legalization under a relatively free-entry regime while some neighboring states remained nonlegal (Texas, Oklahoma and Utah). Firms therefore make their decisions in an environment where cross-border demand is already present and the entry is free. This setting makes it possible to examine how border markets develop when entry is relatively frictionless and how the border premium is reflected in retailer outcomes, rather than asking how incumbent firms respond when a border premium disappears.

Additional demand at the border can show up in two main ways. If entry is relatively open, it can attract more dispensaries, increasing both dispensary density and total sales in the local market. If entry is restricted, the extra demand is more likely to be concentrated among existing retailers, raising revenue per dispensary instead. Which one happens depends on how easily new firms can enter the market.

This paper complements \cite{thomas2019license}, who uses Washington's license-quota system to construct a structural counterfactual of unrestricted entry. While her analysis asks how the market would have developed in the absence of licensing constraints, this paper observes a relatively free-entry market directly. More importantly, it studies how entry responds to a specific source of additional demand: consumers from neighboring states where recreational cannabis remains unavailable. This setting makes it possible to examine whether the border premium is reflected in higher revenue for existing dispensaries or instead attracts additional firms and expands the size of the local market. The results therefore provide empirical evidence on how cross-border demand is distributed across firms when entry is relatively unconstrained.


Using individual dispensary data from New Mexico, including each store’s entry date and location, this paper examines two questions. First, is there evidence of a border premium in cannabis markets near nonlegal-state borders? Second, if such a premium exists, how does it show up in retailer outcomes when entry is relatively unrestricted?


The results show a large and statistically robust border premium at the market level. Border ZIP codes generate about 89 percent more dispensary revenue per resident and have 46 percent greater dispensary density than interior ZIP codes. These effects are larger than the retailer-level response documented by \cite{HansenMillerWeber2020}, who find that sales at Washington retailers near the Oregon border fell by about 36 percent after Oregon opened its recreational cannabis market. One possible explanation is that the two markets differ in how easily supply can respond to cross-border demand. Washington operated under binding retail-license quotas, while New Mexico allows relatively open entry. Consistent with this distinction, \cite{thomas2019license} estimates that removing Washington's license quotas would increase the number of entrants by about 45 percent. The results provide complementary evidence from an actual relatively free-entry market: stronger demand near nonlegal-state borders is associated not only with higher market revenue but also with substantially greater dispensary density. This suggests that when entry is less restricted, part of the border premium may be absorbed through market expansion rather than remaining concentrated in the sales of existing retailers.


The analysis then turns to dispensaries' outcomes within the local ZIP codes. At the firm level, there is no statistically significant difference in either revenue per resident or revenue per store. This strengthens the ZIP-level interpretation because it shows that the larger border market does not translate into a persistent advantage for individual retailers. Instead, the border premium creates greater competition within those markets and spreads the additional demand across more firms.

Finally, this paper asks what happens when border demand meets a more restrictive entry regime. Using county-level data from Colorado, where entry barriers were in place, it constructs a policy counterfactual to examine how those restrictions affect the ability of border markets to capture the cross-border demand. The estimates suggest that the border premium in dispensary density is smaller in Colorado than in New Mexico, though the difference is only marginally significant even when the two states are compared over the same years. The direction is broadly consistent with \cite{thomas2019license}, who shows that license restrictions can substantially limit market expansion. Read as an illustrative benchmark, the estimates imply on the order of \$500 million in foregone Colorado dispensary revenue in 2024, which gives a rough sense of the revenue at stake when entry is restricted in markets exposed to cross-border demand.

The remainder of the paper is organized as follows. Section \ref{litreview} reviews the two strands of literature to which this paper contributes. Section \ref{sec:data} describes the data and institutional setting. Section \ref{sec:empirical} presents the empirical approach used to estimate the border premium and study entry behavior. Section \ref{sec:result1} reports the main results for New Mexico. Section \ref{sec:falsification} presents robustness and sensitivity checks. Section \ref{sec:discussion} discusses what these findings imply about free entry and asks, using a cross-state comparison with Colorado, whether the way border demand is absorbed depends on how entry is regulated. Section \ref{conclusion} concludes.

\section{Related Literature}
\label{litreview}

This paper sits at the intersection of three literatures. The first shows that differences in taxes, prices, and regulations across jurisdictions can redirect consumer demand across borders. The second documents that differences in cannabis legalization generate similar spillovers and that local access to legal retailers matters for cannabis use. The third studies how entry restrictions shape the number of firms, competition, and market outcomes. This paper brings these literatures together by asking how cross-border demand shapes legal cannabis markets and whether its effects depend on how easily new firms can enter.


A large public-finance literature shows that differences in prices, taxes, and access across jurisdictions can lead consumers to shop across borders.
Theoretical work establishes how these shopping opportunities can shape tax competition across jurisdictions \cite{KanburKeen1993}. Empirically, \cite{AsplundFribergWilander2007} show that Swedish alcohol sales respond to prices in neighboring countries and that the response falls sharply with distance from the border. \cite{Lovenheim2008} similarly finds substantial cross-border cigarette purchases in response to price differences, with the response weakening as consumers live farther from a lower-price border. Using Canada--U.S. travel data, \cite{ChandraHeadTappata2014} show more generally that consumers trade off the benefits of purchasing across a border against the cost of traveling there.

Particularly relevant here, cross-border demand can affect not only where consumers purchase but also where retail activity develops. \cite{ManuszakMoul2009} study gasoline markets around tax borders in the Chicago area and show that tax differences substantially shift retail capacity across space. Their estimates imply that the higher-tax side of the border had considerably less gasoline retail capacity than it would have had under equal taxes. This provides an important starting point: persistent differences in demand across borders may change the geography of supply itself. This paper extends this idea to a regulated market in which legal access, rather than only price or taxation, differs across state lines, and studies whether the resulting demand advantage attracts additional firms.

A smaller literature studies cross-border spillovers created specifically by differences in cannabis legalization. \cite{HaoCowan2020} find that recreational legalization in Colorado and Washington was followed by higher marijuana-possession arrests in nearby counties of neighboring nonlegal states and by evidence of increased marijuana use outside the legalizing states, although they caution that some of the arrest response may reflect changes in enforcement. The most closely related paper is \cite{HansenMillerWeber2020}. Exploiting the opening of Oregon's recreational market in October 2015, they find that sales at Washington retailers near the Oregon border fell by approximately 36 percent once Oregon consumers gained legal access at home. Their findings provide direct evidence that differences in legal status across neighboring states can create substantial cross-border retail demand.

This paper starts from this evidence but asks a different question. In \cite{HansenMillerWeber2020}, the relevant Washington retailers were already operating when Oregon opened its market, so the decline in cross-border demand was observed through sales at incumbent stores. This paper studies a setting in which firms make entry and location decisions while the neighboring nonlegal markets are still present. This makes it possible to ask whether cross-border demand is reflected mainly in higher revenue at individual dispensaries or whether it also changes the number of firms the local market supports.

This analysis is also related to work on geographic access to legal cannabis retailers. \cite{AmbroseCowanRosenman2021}, for example, show that adults in Washington use cannabis more often as legal retailers become easier to reach, highlighting the importance of local retail access even after statewide legalization. This literature generally treats proximity and retailer density as determinants of consumer behavior. This paper examines an earlier step in that process: how the spatial distribution of retailers itself responds to local demand. In particular, it asks whether proximity to consumers in neighboring nonlegal states leads to greater dispensary entry and a larger legal market near the border.

The supply response connects this paper to the industrial-organization literature on entry and market structure. The central result in this literature is that larger markets can support more firms and that the entry of additional competitors can change competitive conduct \cite{BresnahanReiss1991, BerryReiss2007}. These forces are especially relevant when regulation determines how easily the number of firms can respond to changes in demand.

Cannabis markets provide unusually clear variation in such regulation. \cite{thomas2019license} develops a structural entry model of Washington's license-quota system and simulates the removal of those restrictions. Her free-entry counterfactual produces substantially more entry and raises total surplus by approximately 18 percent, with much of the welfare loss under the observed system coming from the geographic allocation of licenses. Her analysis shows how license quotas can prevent retail capacity from moving toward locations where it is most valuable.

Related work shows that these restrictions also shape competition among the firms that do enter. Using detailed administrative data from Washington, \cite{HollenbeckUetake2021} find that retail entry restrictions create substantial local market power and show that this market power affects prices, tax incidence, consumer welfare, and government revenue. More recent evidence from Washington's license lotteries reaches a similar conclusion from quasi-random variation in local competition. \cite{HollenbeckHristakevaUetake2026} find that additional nearby competitors reduce retail markups, with much of the competitive effect occurring with the first few entrants. Together, these studies show that entry rules influence both the number of retailers serving a market and how the gains from local demand are distributed among them.

The contribution of this paper is to connect this entry literature directly to cross-border demand. Existing studies establish that differences in cannabis legalization can generate substantial demand across state lines, while a separate literature shows that licensing restrictions influence entry and competition within legal cannabis markets. This paper studies how these two forces interact. New Mexico's relatively open licensing environment makes it possible to observe whether persistent demand from neighboring nonlegal states is associated with greater entry and a larger border market. By comparing ZIP-level market outcomes with individual-dispensary outcomes, this paper examines whether the border premium is concentrated in revenue per firm or instead appears through market expansion and greater competition among retailers. It then compares this pattern against a more restrictive entry environment to provide suggestive evidence on whether entry regulation limits the extent to which border markets expand in response to cross-border demand. The paper therefore moves beyond documenting the existence of cross-border cannabis demand and asks how the structure of the legal market determines where the resulting economic activity appears.


\section{Data and Background}
\label{sec:data}

The empirical setting combines two legal cannabis markets with different approaches to retail entry. New Mexico is the main setting for the analysis. Recreational cannabis sales began in April 2022, and state law does not impose a numerical cap on retail licenses. Licensing is largely administrative, which gives the number and location of dispensaries room to respond to differences in local demand. Colorado provides a useful comparison. Although Colorado does not impose a statewide numerical cap, local governments can restrict or prohibit cannabis businesses, making entry more constrained in many parts of the state. This contrast makes it possible to first study how border markets develop when entry is relatively open and then examine whether the border premium looks different where entry faces greater local restrictions.

The analysis is conducted at three geographic levels. The main border-exposure analysis uses a ZIP-code-by-quarter panel for New Mexico covering all 368 populated ZIP codes from 2022Q2 through 2026Q3, for a total of 6,624 ZIP-quarter observations. To study outcomes at the individual retailer level, this paper uses a dispensary-by-quarter panel covering 1,192 geocoded New Mexico dispensaries over the same period, yielding 10,739 dispensary-quarter observations. Finally, the cross-state comparison uses county-level data from New Mexico and Colorado covering 2022Q2 through 2025Q2 for both states (Section~\ref{sec:crossstate}).

The New Mexico ZIP and dispensary panels are balanced. Because sales are reported monthly at the licensee level, every populated ZIP code is observed in every quarter, and true zero activity can be distinguished from missing observations. This is important for the main outcomes because a ZIP code with no operating dispensary is itself an economically meaningful outcome rather than an observation that should be dropped. The county panel is less complete. Counties enter once dispensary activity is observed and may disappear in periods when revenue is suppressed for confidentiality. As a result, several rural Colorado counties never enter the estimation sample. Appendix Table~\ref{tab:sumstats_state} reports the resulting county coverage: 32 of New Mexico's 33 counties and 42 of Colorado's 64 counties appear in the estimation sample.

Border exposure is defined using geography rather than observed dispensary locations, so that a market's border status is determined independently of firms' entry decisions. A ZIP code, dispensary, or county is classified as border-exposed when two conditions are satisfied. First, its centroid, or the dispensary's own location, must lie within 75 miles of a state line with a neighboring state that does not operate a recreational cannabis market. Second, at least 1,000 residents of that nonlegal state must live within the same 75-mile area. The population requirement excludes sparsely populated stretches of state lines where being physically close to the border does not imply meaningful access to potential cross-border consumers.

For New Mexico, the relevant nonlegal neighbors are Texas, Oklahoma, and Utah. Arizona and Colorado already had operating recreational markets during the sample period and therefore do not provide the same source of captive cross-border demand. Markets close to those legal-state borders are treated as interior for the purposes of the main analysis. Under the baseline definition, 92 of New Mexico's 368 populated ZIP codes are classified as border-exposed, while 276 are classified as interior. At the dispensary level, 436 of the 1,192 geocoded New Mexico dispensaries are border-exposed. At the county level, 10 New Mexico counties and 23 Colorado counties satisfy the corresponding border-exposure definition.

Figure~\ref{fig:spatial_inequality} compares dispensary density and revenue per resident across border and interior ZIP codes in New Mexico; Appendix Figure~\ref{fig:setting} places New Mexico alongside its neighbor Colorado. Border and interior ZIP codes have very similar average populations, roughly 5,700 and 5,600 residents, respectively. This similarity is useful because it means that the two groups are broadly comparable in local population size even though their exposure to neighboring nonlegal-state demand differs substantially.

\begin{figure}[htbp]
\centering
\includegraphics[width=\textwidth]{fig_nm_spatial_inequality.png}
\caption{Spatial inequality in dispensary access, New Mexico. Panel (a): each
bubble is a ZIP-code centroid, sized by average dispensaries per 1,000
residents over 2022Q2--2026Q3 and colored by border status. Panel (b):
population-weighted dispensaries per 1,000 residents (solid bars, left axis) and
revenue per resident (hatched bars, right axis) for border and interior ZIP
codes, pooling all ZIP-quarters.}
\label{fig:spatial_inequality}
\end{figure}

New Mexico cannabis sales come from the Cannabis Control Division's Cannabis Reporting Online Portal \citep{NMCCDReporting}, which reports monthly gross sales by individual licensee and municipality. These data are aggregated to the dispensary, ZIP-code, and county levels depending on the analysis. Colorado sales come from the Colorado Department of Revenue's county-level marijuana sales series \citep{COMarijuanaSalesByCounty}. The available monthly data are aggregated to quarterly for the cross-state comparison.

Dispensary locations in New Mexico come from the Cannabis Control Division's public list of licensed cannabis establishments \citep{NMCCDLicensees}. Store addresses are geocoded using the U.S. Census Bureau's batch geocoder, and each dispensary is assigned to a ZIP code and county. Colorado dispensary information comes from the Marijuana Enforcement Division's licensee records \citep{COMedLicensees}, which are also geocoded and assigned to counties. Because New Mexico reports sales at the licensee level, operating dispensaries can be identified quarter by quarter. Historical Colorado counts are less complete and must be reconstructed from the available license records.

Population and geographic data come from several sources. ZIP-code populations and centroids are taken from the SimpleMaps U.S. ZIP database \citep{SimpleMapsZIP}, which is based on American Community Survey data. County population comes from the U.S. Census Bureau's Population Estimates Program, using the 2023 vintage \citep{CensusPopEst2023}. Population surrounding individual dispensaries is calculated using 2020 Census block populations and boundaries \citep{CensusBlocks2020}. State and county boundaries come from the Census Bureau's 2023 cartographic files \citep{CensusCountyShapefile}. For the falsification exercises in Section~\ref{sec:falsification}, this paper also uses non-cannabis retail data from the New Mexico Taxation and Revenue Department's gross-receipts reports \citep{NMTRDGrossReceipts} and the Census Bureau's ZIP Codes Business Patterns \citep{CensusZBP}.

This paper focuses on three main outcomes. Revenue per capita is total cannabis sales in a market divided by its resident population. Dispensary density is the number of operating dispensaries per 1,000 residents. Revenue per dispensary is total cannabis sales divided by the number of operating dispensaries and is therefore defined only for markets with at least one store. The first two outcomes describe the overall size of the local cannabis market, while revenue per dispensary captures the amount of sales generated by the average retailer. Where outcomes enter the regression in logs, coefficients are translated into percentage differences using $(e^{\hat\beta}-1)\times100$. Table~\ref{tab:sumstats_nm} summarizes the main New Mexico variables, and Appendix Table~\ref{tab:sumstats_state} the cross-state county sample. 



\begin{table}[htbp]
\centering
\caption{Summary statistics, New Mexico}
\label{tab:sumstats_nm}
\footnotesize
\begin{tabular}{lcccc}
\toprule
 & \multicolumn{2}{c}{Full sample} & Border & Interior \\
\cmidrule(lr){2-3}\cmidrule(lr){4-4}\cmidrule(lr){5-5}
Variable & Mean & SD & Mean & Mean \\
\midrule
\multicolumn{5}{l}{\textit{Panel A: ZIP-code level (quarterly, 2022Q2--2026Q3)}} \\
Revenue per capita (\$/quarter) & 130.5 & 769.1 & 354.6 & 55.7 \\
Dispensaries per 1,000 residents & 0.819 & 4.890 & 1.732 & 0.514 \\
Revenue per dispensary (\$000s/quarter) & 182.2 & 159.8 & 195.5 & 175.8 \\
ZIP population & 5,588 & 11,709 & 5,688 & 5,555 \\
Distance to nonlegal state line (mi) & 106.8 & 56.4 & 29.7 & 132.6 \\
Nonlegal population within 75 mi (000s) & 72.1 & 205.5 & 288.2 & 0.0 \\
\addlinespace
ZIP codes & 368 & & 92 & 276 \\
ZIP--quarters & 6,624 & & 1,656 & 4,968 \\
\midrule
\multicolumn{5}{l}{\textit{Panel B: Dispensary level (quarterly, 2022Q2--2026Q3)}} \\
Revenue (\$000s/quarter) & 208.8 & 317.8 & 221.8 & 201.8 \\
Population within 1 km & 3,101 & 2,183 & 2,340 & 3,541 \\
Population within 25 mi (000s) & 380.6 & 341.5 & 206.9 & 480.8 \\
Competing dispensaries within 25 mi & 100 & 95 & 37 & 135 \\
Dispensary age (quarters) & 5.7 & 4.5 & 5.4 & 5.8 \\
Distance to nonlegal state line (mi) & 109.4 & 76.3 & 21.5 & 160.1 \\
\addlinespace
Dispensaries & 1,192 & & 436 & 756 \\
Dispensary--quarters & 10,739 & & 3,806 & 6,933 \\
Distinct brands & 684 & & & \\
\bottomrule
\end{tabular}
\begin{minipage}{0.92\textwidth}
\vspace{0.4em}\footnotesize
\textit{Notes:} New Mexico, 2022Q2 through 2026Q3. A ZIP code (Panel A) and dispensary (Panel B) is classified as \emph{border} if its centroid lies within 75 miles of the New Mexico line with a state that has no operating recreational cannabis market (Texas, Oklahoma, or Utah) and at least 1{,}000 residents of that nonlegal state live within the same 75-mile reach; Arizona and Colorado are legal-market neighbors and count as interior. Panel B local-resident population is the 2020 Census block population within a 1-km radius of each dispensary.
\end{minipage}
\end{table}


The data also have several limitations. First, the New Mexico panels begin when recreational sales start in April 2022. The analysis therefore compares border and interior markets after legalization rather than estimating the causal effect of legalization itself. Second, population is measured using a fixed vintage and is not adjusted for the share of residents who are legally old enough to purchase cannabis. Third, Colorado's historical dispensary counts are reconstructed from license records rather than observed from a complete quarter-by-quarter operating panel. This introduces greater measurement error into Colorado's density and revenue-per-dispensary outcomes, although it does not affect county-level revenue per capita. Finally, this paper does not observe the residence of individual cannabis consumers. Cross-border demand is instead measured using the size and proximity of the population living in neighboring nonlegal states. The empirical strategy therefore identifies differences associated with exposure to potential cross-border demand rather than directly observing individual cross-border purchases.

\section{Empirical Strategy}
\label{sec:empirical}

The estimand is the effect of exposure to
cross-border demand from non-legalized neighboring states on outcomes in
New Mexico's dispensary market. Border premium can be translated as any
resulting difference in market outcomes between exposed (border) and unexposed (interior) markets. The idea is that legalization creates the legal market suddenly. Given that timing was not fully anticipated by firms or consumers, then the policy provides a relatively sharp market-opening shock. After the market opens, ZIP codes near nonlegal states are suddenly able to serve an additional pool of consumers that interior ZIP codes cannot serve as easily.

\subsection{Border-Exposure Panel Design}
\label{sec:borderexposure}

The main empirical analysis compares cannabis-market outcomes in New Mexico ZIP codes that are exposed to potential demand from neighboring states without legal recreational cannabis markets. Because the sample begins after recreational sales started in New Mexico, the analysis is not designed to estimate the effect of legalization itself. Instead, it asks whether legal cannabis markets located close to nonlegal neighboring states are systematically larger than comparable interior markets.

The analysis begins with a static specification that estimates the average border premium over the full sample period:

\begin{equation}
\label{baseline}
\log(y_{zt})=\alpha
+
\beta \text{Border}_z
+
\gamma \log(\text{Population}_z)
+
\delta_t
+
\varepsilon_{zt}.
\end{equation}

The outcome \(y_{zt}\) is one of three cannabis-market outcomes in ZIP code \(z\) and quarter \(t\), each in natural logarithms: total dispensary revenue, the number of operating dispensaries, and revenue per operating dispensary. The coefficient of interest, \(\beta\), is the border premium -- the average difference in the log outcome between border and interior ZIP codes, holding local population fixed. For the first two outcomes this is equivalent to a per-resident specification (revenue per capita, dispensaries per capita); writing the outcome in levels keeps \(\log(\text{Population}_z)\) as a right-hand-side control rather than having population appear on both sides of the equation.

\(\text{Population}_z\) is the resident population of ZIP code \(z\) itself. Its coefficient \(\gamma\) captures how each outcome scales with the size of the local, legal-side market; it is not a measure of cross-border or captive demand, which enters the specification only through \(\text{Border}_z\). Quarter fixed effects, \(\delta_t\), absorb factors that affect the entire New Mexico cannabis market at the same time, including seasonality, statewide market growth, and common price movements. Standard errors are clustered at the county level to allow outcomes within the same county to be correlated over time.

Equation \ref{baseline} provides a single estimate of the average border premium. It is complemented with a more flexible specification that allows the border premium to differ across quarters:

\begin{equation}
\label{baseline_dynamic}
\log(y_{zt})=
\alpha
+
\sum_q
\beta_q
\left(
\text{Border}_z
\times
\mathbbm{1}[\text{Quarter}=q]
\right)
+
\gamma \log(\text{Population}_z)
+
\delta_q
+
\varepsilon_{zt}.
\end{equation}

In this specification, each \(\beta_q\) directly measures the difference between border and interior ZIP codes in quarter \(q\). This avoids relying on a single reference quarter and makes the estimates easier to interpret. The dynamic specification shows whether the border premium is persistent, grows, or declines as New Mexico's recreational cannabis market develops.

Border status is determined entirely by geography and does not depend on where dispensaries choose to locate. The ZIP code is classified as border-exposed ($Border_z=1$) if its centroid lies within 75 miles of a New Mexico border with Texas, Oklahoma, or Utah and at least 1,000 residents of the neighboring nonlegal state live within the same 75-mile market area. The population requirement of 1,000 residents excludes sparsely populated stretches of the state line where there is no meaningful potential for cross-border demand (like unpopulated desert areas).
Under the geography-based definition of border status, 92 of New Mexico's 368 populated ZIP codes are classified as border-exposed and 276 as interiors.

Dispensaries close to Arizona and Colorado borders, two neighboring states with New Mexico, are considered as interior because recreational cannabis was already legally available in those states during the sample period.

The estimation sample is the New Mexico ZIP codes that keeps all 368 populated ZIP codes, including those that never record a dispensary at some point over 2022Q2--2026Q3. 

The identification assumption requires that comparable border and interior ZIP codes would have had similar cannabis-market outcomes in the absence of cross border demand. This assumption could fail if border areas are systematically different for reasons unrelated to cannabis. For example, border ZIP codes could have greater commercial activity, different transportation access, or other characteristics that independently attract retailers. We run several placebo checks to ensure the robustness of our results.

First, non-cannabis retail sectors are examined to determine whether border areas are simply more commercially active in general. Second, ZIP codes near nonlegal-state borders are compared with ZIP codes near Arizona and Colorado, where recreational cannabis was already legal. This provides a useful placebo test for whether border areas are systematically different from interior areas for reasons unrelated to cross-border cannabis demand.

 Third, a continuous measure of distance replaces the variable $Border$ in specification \ref{baseline_dynamic}. The argument is that, if proximity to consumers in neighboring nonlegal states is driving the effect, the border premium should increase gradually with greater exposure, rather than appearing only because of a specific cutoff. Forth, the study uses an alternative control group consisting of ZIP codes along New Mexico’s borders with already-legal states, Colorado and Arizona. If border regions are systematically different from interior areas simply because they are near a state border, we should not observe a significant difference between ZIP codes bordering nonlegal states and those bordering already-legal states.
 Finally, 50- and 100-mile distance thresholds are used to define whether ZIP codes are classified as border areas. This allows a broader sample of ZIP codes with more diverse characteristics These exercises do not prove the identifying assumption, but they help distinguish the cross-border cannabis-demand mechanism from several important alternative explanations. Results from these tests are reported in Section~\ref{sec:falsification}.

\section{Results}
\label{sec:result1}


\subsection{The Border Premium in New Mexico}

Table~\ref{tab:nm_baseline} reports the baseline border-exposure estimates (equation~\ref{baseline}).Dispensary revenue per resident is 89.4 percent higher in border ZIP codes than in interiors \(\left(e^{0.639}-1\right)\).
Border ZIP codes also support substantially more dispensaries per resident (+46.7 percent), and 29.1 percent more revenue per
dispensary. 

Figure~\ref{fig:dynamic_borderpremium} plots the dynamic border-premium estimation results (equation~\ref{baseline_dynamic}).
These quarter-specific coefficients are not read as pre-trend or post-treatment effects, but as a description of whether the border-interior gap is persistent or evolves as
New Mexico's recreational market develops. Revenue per capita and dispensary density are positive in most
quarters and significant in the majority of them;
revenue per dispensary is noisier and significant in fewer quarters (7 of 18), consistent with the free-entry
mechanism. If an existing store starts earning unusually high revenue, that attracts new stores into the market. As more stores enter, the same customer demand is spread across more businesses, which brings the average revenue per store back down.

Taken together, the results point to two main conclusions. First, New Mexico exhibits a clear border premium in its legal cannabis market. ZIP codes near nonlegal-state borders generate substantially more cannabis revenue per capita and support more dispensaries per capita than interiors. Importantly, these differences persist across most of the post-legalization period rather than being driven by a small number of quarters. This persistent gap is the paper’s main finding and provides the strongest evidence of additional demand associated with proximity to nonlegal states.

Second, the way this additional demand is distributed across retailers is consistent with the free entry market characteristic. Unlike revenue per capita and dispensary density, revenue per dispensary is not persistently higher in border markets, reaching statistical significance in only 7 of 18 quarters. This pattern suggests that higher border demand does not translate into a lasting revenue advantage for individual dispensaries. Instead, temporarily higher revenues attract additional entry, increasing the number of stores over which that demand is shared. As a result, the border premium appears primarily through a larger market rather than through persistently higher revenue for incumbent retailers. 

Compared with the existing literature, the estimates suggest that border areas have substantially more cannabis-market activity overall, reflected in both higher revenue per resident and more dispensaries per resident.

The estimate for revenue per dispensary is similar in magnitude to that reported by \cite{HansenMillerWeber2020}, with dispensaries in New Mexico border ZIP codes earning about 29 percent more $\left[(e^{0.255}-1)*100\right]$ than those in interiors.
The larger effects on revenue per capita and dispensary density suggest that, in New Mexico, cross-border demand is reflected not only in higher sales at existing stores but also in additional entry. This distinction is consistent with the institutional setting: whereas Washington restricted retail entry through licensing limits, New Mexico allows entry to respond more freely to local demand conditions. The results therefore suggest that under freer entry, part of the border premium is absorbed through market expansion rather than remaining concentrated in the revenues of incumbent retailers.


\begin{table}[htbp]
\centering
\caption{Border Premium in New Mexico}
\label{tab:nm_baseline}

\begin{tabular*}{\textwidth}{@{\extracolsep{\fill}}lccc}
\toprule
& (1) & (2) & (3) \\
& \shortstack{Dispensary\\revenue}
& \shortstack{Number of\\dispensaries}
& \shortstack{Revenue per\\dispensary} \\
& $\log(y)$
& $\log(y)$
& $\log(y)$ \\
\midrule

Border ZIP
    & $0.639^{***}$
    & $0.383^{***}$
    & $0.255^{*}$ \\
    & (0.223)
    & (0.120)
    & (0.149) \\

\addlinespace

$\log(\text{Population})$
    & $0.561^{***}$
    & $0.319^{***}$
    & $0.242^{***}$ \\
    & (0.078)
    & (0.046)
    & (0.051) \\

\midrule

Quarter fixed effects
    & Yes & Yes & Yes \\

Observations
    & 1,704
    & 1,704
    & 1,704 \\

Number of ZIP codes
    & 116
    & 116
    & 116 \\

$R^{2}$
    & 0.271
    & 0.257
    & 0.167 \\

\bottomrule
\end{tabular*}

\begin{minipage}{\textwidth}
\vspace{0.4em}
\fontsize{9pt}{9pt}\selectfont

\textit{Notes:}
The unit of observation is a ZIP--quarter. The sample covers New Mexico from
2022Q2 through 2026Q3 and is restricted to the 116 ZIP codes that record at
least one operating dispensary at some point in the sample. The outcome is the
natural log of total quarterly dispensary revenue in column~(1), the log of the
number of operating dispensaries in column~(2), and the log of revenue per
operating dispensary in column~(3).

$\text{Population}$ is the ZIP code's own resident population. Its coefficient is
the elasticity of the outcome with respect to local population and serves as a
scale control; it is not a measure of cross-border or captive demand, which
enters only through Border ZIP. For columns~(1) and~(2) this specification is
algebraically identical to regressing the corresponding per-resident outcome
(revenue per capita, dispensaries per capita) on the same regressors: the Border
coefficient is unchanged and only the population coefficient shifts by one. In
percentage terms ($e^{\hat\beta}-1$), border ZIP codes have about 89 percent
more dispensary revenue, 47 percent more dispensaries, and 29 percent more
revenue per dispensary than interior ZIP codes with the same resident
population.

Border ZIP is an indicator equal to one if the ZIP-code centroid is within
75 miles of the border with Texas, Oklahoma, or Utah and at least 1,000
residents of those nonlegal states live within the same 75-mile radius.
All specifications include quarter fixed effects. Standard errors, reported in
parentheses, are clustered at the county level.
\end{minipage}

\end{table}


\begin{figure}[!ht]
\centering

\makebox[\textwidth][c]{%
    \includegraphics[width=1.25\textwidth]{fig_nm_dynamic_borderpremium.png}
}

\captionsetup{font=ninept}
\caption{Dynamic Border Premium by Quarter, New Mexico.
$\text{Border}_z \times \text{Quarter}_q$ coefficients from
equation~\ref{baseline_dynamic}, with 2022Q2 as the reference quarter.
The figure reports county-clustered 95\% confidence intervals.
The estimates show whether the border-interior difference persists or
changes as the market develops; they should not be interpreted as
pre-trend or post-treatment effects of legalization.}
\label{fig:dynamic_borderpremium}
\end{figure}

\FloatBarrier

\newpage
\subsection{Entry as the Supply-Side Response}
\label{sec:mechanism}

Per-capita revenue is the product of dispensaries per capita and revenue per dispensary, so it can be decomposed exactly as
$\log(\text{RevenuePerCapita}) = \log(\text{DispensariesPerCapita}) + \log(\text{RevenuePerDispensary})$.

This decomposition separates the ZIP-level border premium into an extensive margin, measured by dispensaries per capita, and an intensive margin, measured by revenue per dispensary. The results in Section \ref{sec:result1} show that border ZIP codes have higher revenue per capita and greater dispensary density, while the difference in revenue per dispensary is smaller and less precisely estimated. However, the imprecise ZIP-level estimate for revenue per dispensary alone is not sufficient to establish this mechanism.


The mechanism is therefore examined more directly using dispensary-level data. Each dispensary is classified as border-exposed using the same 75-mile and 1,000-resident criterion used in the ZIP-level analysis, and the following dynamic specification is estimated:

\begin{equation}
\label{dispensary_dynamic}
\log(y_{dt})=
\alpha
+
\sum_q
\beta_q
\left(
\text{Border}_d
\times
\mathbbm{1}[\text{Quarter}=q]
\right)
+
\gamma \log(\text{Population}_d)
+
\delta_q
+
\lambda_{b(d)}
+
e_{dt}
\end{equation}

In this specification $y_{dt}$ is revenue per local resident for dispensary $d$ in quarter $q$, and $\text{Population}_d$ is that same one-kilometre population, kept as a control so that its elasticity need not equal $-1$. The term $\delta_q$ gives quarter fixed effects, and $\lambda_{b(d)}$ is a fixed effect for the retail brand (chain) that operates dispensary $d$; it absorbs chain-level differences in pricing, product assortment, and store format. Figure \ref{fig_dispensary} shows no persistent difference in revenue per local resident between border and interior dispensaries over time. The estimates fluctuate around zero and do not display a systematic upward pattern in border markets. This evidence strengthens the interpretation of the ZIP-level results: The border premium appears primarily through market expansion in border areas, reflected in higher aggregate revenue and greater entry rate.

This pattern is also consistent with the standard industrial-organization logic of free entry \citep{BresnahanReiss1991}:
a border location may generate additional local demand, but when entry is not constrained, that demand attracts new firms. As entry expands, the additional market is shared across a larger number of dispensaries, reducing any persistent revenue margin for individual border retailers. 


\begin{figure}[!ht]
\centering

\makebox[\textwidth][c]{%
    \includegraphics[width=1\textwidth]{fig_eventstudy_nm_dispensary_revperresident_1km_vs_2km.png}
}

\captionsetup{font=ninept}
\caption{Revenue per local resident, one- versus two-kilometre denominator.
Border-by-quarter coefficients from specification ~\ref{dispensary_dynamic}, with the resident-population denominator computed
over a one-kilometre radius (as in the main specification) and a two-kilometre
radius. County-clustered
95 percent confidence intervals.}
\label{fig_dispensary}
\end{figure}

\section{Robustness checks}
\label{sec:falsification}

Border ZIP codes may differ from interior ZIP codes for reasons unrelated to cannabis. More generally, border areas may have higher levels of commercial activity across many sectors, which could make them systematically different from interior areas even in the absence of cross-border cannabis demand. A related concern is that the estimated border premium may depend on the particular cutoff used to define a border ZIP code. For example, a reader may reasonably ask why 75 miles is the appropriate threshold and whether the results would change if a narrower or wider distance were used. The flexible empirical design makes it possible to address these concerns one at a time by changing key features of the baseline specification and examining whether the main results remain robust.
 
First, dispensaries are sorted into 25-mile distance bands, using the band closest to the border (0--25 miles) as the reference group, and the following is re-estimated

\begin{equation}
\log(y_{it}) = \alpha + \sum_{k} \beta_k \cdot \mathbbm{1}[\text{DistanceBand}_i = k] + \gamma\cdot\log(\text{Population}_i) + \delta_t + \varepsilon_{it}
\end{equation}

at the dispensary level. So, the variable $Population_i$ represents the population within 1km radius of where the dispensary is located. As shown in Figure \ref{fig:nm_doseresponse}, revenue per dispensary and revenue per resident remain relatively stable as distance from the nonlegal-state border increases. The only outcome that declines systematically with distance is dispensary density. In particular, there is no statistically significant decline in the number of dispensaries within the first 75 miles from the border; the first clear reduction appears in the 75--100 mile distance band. This pattern helps address the concern that the baseline 75-mile cutoff is driving the results. If the estimated border premium were simply an artifact of where the boundary was drawn, the results would be expected to change sharply when moving across nearby distance bands. Instead, the main pattern remains broadly similar across the distance distribution. The premium also strengthens with the size of the nearby nonlegal-state population (Appendix Table~\ref{tab:nm_nonlegalpop_doseresponse} and Figure~\ref{fig:nm_nonlegalpop_doseresponse}).

Border status is also re-defined using alternative 50-mile and 100-mile thresholds. In the baseline specification, a ZIP code is classified as border-exposed if its centroid lies within 75 miles of a nonlegal-state border and at least 1,000 residents of the neighboring nonlegal state live within the same 75-mile radius. The main specifications are re-estimated after replacing the 75-mile threshold with 50- and 100-mile thresholds. The point estimates are broadly similar to the baseline 75 mile distance across all three definitions \ref{fig:eventstudy_50_100mi}, and confidence intervals become wider under the 100-mile definition, which is expected because this broader cutoff includes ZIP codes that are less directly exposed to the neighboring nonlegal market. Overall, the main findings do not depend on treating 75 miles as cutoff for treating the ZIP codes as border. 

The identification strategy relies on the assumption that border and interior regions would have followed similar patterns in the absence of cross-border demand from nonlegal states. One concern is that the estimated border premium may reflect factors related to being near a state border, rather than demand from consumers in neighboring states where cannabis is not legal.
 The baseline uses interior New Mexico ZIP codes as the comparison group, and these may differ from border ZIP codes for reasons that have nothing to do with cannabis. To separate the two explanations, only ZIP codes that sit near a state line are kept, split by the kind of neighbor they face. The treated group is unchanged: border ZIP codes within 75 miles of the Texas, Oklahoma, or Utah line, with at least 1{,}000 residents of that nonlegal state within reach. The new control group is New Mexico ZIP codes that meet the same 75-mile and 1{,}000-resident rule for the Colorado or Arizona line, where consumers on the other side can already buy cannabis legally at home. If the premium comes from cross-border demand it should persist in this comparison; if it comes from a generic border effect it should disappear, since both groups are now close to a border. Figure~\ref{fig:dynamic_legalcontrol} shows that it persists. Border ZIP codes keep a large and precisely estimated premium in revenue per capita and in dispensary density relative to the legal-border control group. Revenue per dispensary again shows no border gap, consistent with New Mexico's free entry channeling the extra demand into more stores rather than higher revenue per store.

The paper also tests whether border areas are generally more commercially active than interior areas, regardless of cross-border cannabis demand.
 The falsification tests use two sectors with similar cross-state policies: alcohol retailers and grocery establishments. Using ZIP-quarter data, they ask whether the border-interior difference in these sectors changes when recreational cannabis becomes legal in New Mexico. The relevant test is therefore not whether border and interior ZIP codes ever differ in alcohol or grocery activity, since such differences may exist for many unrelated reasons, but whether those differences change around the legalization of recreational cannabis.

To examine this, border status is interacted with each calendar quarter and the dynamic estimates are anchored around New Mexico's legalization date. Figure~\ref{fig:noncannabis_eventstudy} plots the resulting border coefficients over time. The key question is whether activity in these unrelated sectors shows a similar increase around legalization date or not. It does not, the border-interior difference in alcohol retail activity shows no sharp change around legalization; any pre-existing difference simply continues afterward. Grocery establishments show little evidence of a systematic difference between border and interior ZIP codes over the sample period. This suggests that the increase in cannabis activity near the border is not simply part of a broader expansion in commercial activity in the same areas. 
Taken together, these robustness checks provide evidence that the main results are not particularly sensitive to the potential confounders considered above. However, they cannot completely rule out all possible sources of confounding, nor is doing so the primary scope of this paper.



\section{Discussion}
\label{sec:discussion}

The New Mexico evidence shows how a relatively free-entry cannabis market absorbs cross-border demand. Legal cannabis markets near the border are larger than in the state's interior, a difference referred to as the border premium throughout the paper: border ZIP codes have 89.4 percent higher dispensary revenue per resident and 46.7 percent more dispensaries per resident than interior ZIP codes. The premium appears primarily through greater dispensary density rather than through significantly higher revenue per store. This is consistent with the standard free-entry prediction that extra demand attracts more firms rather than raising the revenue of existing ones \citep{BresnahanReiss1991}.

This paper next asks whether the size of the border premium depends on how market entry is regulated. To study this, it compares two policy environments -- Colorado, where local governments can restrict entry, and New Mexico, where entry is largely unrestricted -- and compares dispensary revenue in border counties across the two settings. This comparison also shows what policymakers may give up when they restrict entry into sin-good markets. Measuring the revenue lost under entry restrictions provides a simple measure of the economic cost of these policies. This can help policymakers compare these costs with the goals the restrictions are meant to achieve.

Appendix Table~\ref{tab:regime_interaction} reports the $\text{Border}_c \times \text{Restricted}_s$ interaction for all three outcomes; Appendix Section~\ref{sec:crossstate} describes the county-level sample and specification. The interaction is negative for each one, so the border premium looks smaller in Colorado than in New Mexico, but it is estimated imprecisely. It is marginally significant for dispensary density ($\theta=-0.932$) and not statistically distinguishable from zero for revenue per capita ($\theta=-0.458$) or revenue per store. The direction of the difference still lines up with how the two markets are organized. In Colorado, where entry is restricted, the border premium shows up mainly in higher revenue per store, an intensive-margin response. In New Mexico, where entry is largely unrestricted, it shows up mainly in a larger number of stores, an extensive-margin response.

To put the interaction in dollar terms, this paper uses the dispensary-density estimate, the outcome for which the difference across entry regimes is statistically significant. It asks how much more dispensary revenue Colorado's border counties would have generated in 2024 if their density border premium had matched New Mexico's rather than the restricted-entry premium, holding revenue per store at its observed Colorado level. Applying the multiplier $e^{-\theta}\approx2.54$ (using $\theta=-0.932$) to each border county's own 2024 revenue and summing, Colorado's border counties earned \$324.2 million in 2024, while the free-entry benchmark puts them at \$823 million, a gap of roughly \$499 million (Appendix Table~\ref{tab:counterfactual}). The interaction is only marginally significant and the bootstrap interval for this gap still reaches zero, so the dollar figure is read as an illustrative benchmark rather than a precise estimate of foregone revenue.

Because the entry regime is not randomly assigned and only two states are compared directly, these cross-state estimates are treated as suggestive evidence that the border premium differs across entry regimes, rather than as a clean causal estimate of the effect of licensing restrictions. Read this way, the comparison reinforces rather than competes with the New Mexico findings. The direction is consistent with \citet{thomas2019license}, who shows that license restrictions can substantially limit market expansion, and with evidence that entry restrictions in cannabis markets create local market power for incumbent retailers \citep{HollenbeckUetake2021}.

Taken together, the results suggest that whether cross-border demand expands a legal cannabis market or mainly raises revenue for existing stores depends on state entry regulations. For policymakers weighing licensing and zoning rules, this is the margin at stake: tighter entry restrictions may hold down the number of dispensaries near a nonlegal-state border, shifting the adjustment toward higher revenue at incumbent stores. The findings also give policymakers a tool to weigh higher revenue against the social costs of expanding sin-good markets.

\section{Conclusion}
\label{conclusion}

This study analyzes dispensary outcomes in U.S. cannabis markets where different entry policies are implemented
in different states. It first considers a free-entry market in New Mexico, where legal cannabis markets near the border are larger than in the state’s interior, a difference referred to as the border premium throughout the paper. Border ZIP codes, which are more exposed to demand from neighboring nonlegal states, have 89.4 percent higher dispensary revenue and 46.7 percent more dispensaries per resident.
The premium appears
primarily through greater dispensary density rather than through significantly higher revenue per store. This is consistent with the standard free-entry prediction that extra demand attracts more firms rather than raising the revenue of existing ones.

Using Colorado, where entry into the legal cannabis market is more restricted, this study compared how the border premium appears under a different regulatory regime. The results suggest that in Colorado, the border premium shows up mainly through higher revenue per dispensary rather than through a larger number of dispensaries in border markets. This is interpreted as suggestive evidence that entry restrictions can limit how much border cannabis markets expand in response to cross-border demand.

Taken together, the results suggest that the extent to which cross-border demand increases dispensary revenue in border markets depends on state entry regulations.
The findings also give policymakers a tool to weigh higher revenue against the social costs of expanding sin-good markets.

\newpage

\appendix
\section{Appendix}
\renewcommand{\thetable}{A.\arabic{table}}
\renewcommand{\thefigure}{A.\arabic{figure}}
\renewcommand{\theequation}{A.\arabic{equation}}
\setcounter{table}{0}
\setcounter{figure}{0}
\setcounter{equation}{0}

\subsection{Additional Results for New Mexico}
\label{app:nm}


\begin{table}[htbp]
\centering
\caption{New Mexico dispensary-level border effect, noise-minimizing specification (grid-search selected)}
\label{tab:nm_dispensary_border_best}
\footnotesize
\begin{tabular}{lccccc}
\toprule
Outcome & Coef. (\%) & SE & $p$ (cl.) & $n$ & Brand FE \\
\midrule
Revenue per capita & $+0.149$ ($+16.1\%$) & 0.107 & 0.162 & 10725 & Yes \\
Revenue (per dispensary) & $+0.150$ ($+16.2\%$) & 0.103 & 0.144 & 10739 & Yes \\
Local dispensary density & $-0.154$ ($-14.3\%$) & 0.160 & 0.336 & 10739 & No \\
\bottomrule
\end{tabular}
\begin{minipage}{0.98\textwidth}
\vspace{0.3em}
\footnotesize
\textit{Notes:} Sample: geocoded New Mexico dispensaries with valid coordinates (1,192 of 1,253). Border =
dispensary within 75 miles of the nonlegal-state line, with $\geq$1,000 nonlegal-state residents within that
reach. Population control (in logs), quarter fixed effects, county-clustered standard errors throughout;
chain/brand fixed effects included for the two revenue outcomes only (see text). Compare to the naive
baseline specification (county population control, no brand FE, unweighted, county-clustered): standard
errors there were 0.129, 0.125, and 0.272 respectively -- 17, 18, and 41 percent higher than the
noise-minimizing specification reported here.
\end{minipage}
\end{table}

\begin{table}[htbp]
\centering
\caption{Dose-response of dispensary outcomes to distance from the nonlegal-state line, New Mexico}
\label{tab:nm_doseresponse}
\footnotesize
\begin{tabular}{lccc}
\toprule
Distance band (mi) & Revenue per capita & Revenue per dispensary & Dispensary density \\
\midrule
0-25 & ref. & ref. & ref. \\
25-50 & $+0.275$ ($+31.6\%$), $p=0.244$ & $+0.277$ ($+31.9\%$), $p=0.235$ & $+0.083$ ($+8.6\%$), $p=0.827$ \\
50-75 & $+0.054$ ($+5.5\%$), $p=0.810$ & $+0.044$ ($+4.5\%$), $p=0.838$ & $-0.217$ ($-19.5\%$), $p=0.276$ \\
75-100 & $+0.166$ ($+18.0\%$), $p=0.432$ & $+0.150$ ($+16.2\%$), $p=0.476$ & $-0.634$ ($-46.9\%$), $p=0.002$ \\
100-125 & $+0.289$ ($+33.5\%$), $p=0.315$ & $+0.291$ ($+33.8\%$), $p=0.308$ & $-0.473$ ($-37.7\%$), $p=0.095$ \\
125-150 & $+0.257$ ($+29.3\%$), $p=0.244$ & $+0.259$ ($+29.6\%$), $p=0.233$ & $-0.536$ ($-41.5\%$), $p=0.002$ \\
150-175 & $+0.116$ ($+12.3\%$), $p=0.561$ & $+0.109$ ($+11.5\%$), $p=0.580$ & $-0.675$ ($-49.1\%$), $p<0.001$ \\
\bottomrule
\end{tabular}
\begin{minipage}{0.98\textwidth}
\vspace{0.3em}
\footnotesize
\textit{Notes:} Each column is a separate regression of the named log outcome on the full set of distance-band
indicators (omitted reference: 0-25 miles), a log population control, and quarter fixed effects, New Mexico
only. The 175-200 mile band is excluded from the sample (see text). Revenue outcomes: $n=6{,}682$-$6{,}686$
dispensary-quarters. Dispensary density: $n=1{,}302$ ZIP-quarters. County-clustered standard errors.
\end{minipage}
\end{table}


\begin{table}[htbp]
\centering
\caption{Dose-response of the border premium to nearby nonlegal-state population, New Mexico}
\label{tab:nm_nonlegalpop_doseresponse}
\footnotesize
\begin{tabular}{lcccc}
\toprule
 & \multicolumn{2}{c}{Revenue per capita} & \multicolumn{2}{c}{Dispensary density} \\
\cmidrule(lr){2-3}\cmidrule(lr){4-5}
Nonlegal pop.\ within 75 mi (000s) & Coef. & SE & Coef. & SE \\
\midrule
0 (reference) & ref. & & ref. & \\
0--50 & $-0.019$ & (0.221) & $+0.187$ & (0.138) \\
50--250 & $+0.869^{**}$ & (0.356) & $+0.465^{***}$ & (0.179) \\
250+ & $+0.859^{***}$ & (0.238) & $+0.437^{***}$ & (0.150) \\
\midrule
Quarter fixed effects & Yes & & Yes & \\
Observations & 1,704 & & 1,704 & \\
Number of ZIP codes & 116 & & 116 & \\
\bottomrule
\end{tabular}
\begin{minipage}{0.92\textwidth}
\vspace{0.4em}\footnotesize
\textit{Notes:} Each column reports $\beta_k$ from $\log(y_{zt}) = \alpha + \sum_k \beta_k\,\mathbbm{1}[\text{bin}_z=k] + \gamma\log(\text{Population}_z) + \delta_t + \varepsilon_{zt}$, estimated on the 116 New Mexico ZIP codes with any dispensary activity. $\text{Population}_z$ is the ZIP's own resident population. Quarter fixed effects, county-clustered standard errors. $^{*}p<0.10$, $^{**}p<0.05$, $^{***}p<0.01$.
\end{minipage}
\end{table}


\begin{table}[htbp]
\centering
\caption{Border premium including ZIP codes with no dispensary activity}
\label{tab:nm_baseline_zeroinclusive}

\begin{tabular*}{\textwidth}{@{\extracolsep{\fill}}lccc}
\toprule
& (1) & (2) & (3) \\
& \shortstack{Revenue\\per capita}
& \shortstack{Dispensary\\density}
& \shortstack{Revenue\\per dispensary} \\
& $\operatorname{asinh}(y)$
& $\operatorname{asinh}(y)$
& $\log(y)$ \\
\midrule

Border ZIP
    & $0.874^{***}$
    & $0.257^{***}$
    & $0.255^{*}$ \\
    & (0.226)
    & (0.075)
    & (0.149) \\

\addlinespace

$\log(\text{Population})$
    & $0.448^{***}$
    & 0.012
    & $0.242^{***}$ \\
    & (0.083)
    & (0.020)
    & (0.051) \\

\midrule

Quarter fixed effects
    & Yes & Yes & Yes \\

Observations
    & 6,624
    & 6,624
    & 1,704 \\

Number of ZIP codes
    & 368
    & 368
    & 116 \\

$R^{2}$
    & 0.146
    & 0.034
    & 0.167 \\

\bottomrule
\end{tabular*}

\begin{minipage}{\textwidth}
\vspace{0.4em}
\fontsize{9pt}{9pt}\selectfont

\textit{Notes:} Companion to Table~\ref{tab:nm_baseline}. Columns~(1) and~(2)
apply the inverse hyperbolic sine transformation to revenue per capita and to
dispensaries per 1{,}000 residents, so that the 368 populated ZIP codes with no
recorded dispensary activity remain in the sample; the border premium is close
to the log-level estimates in Table~\ref{tab:nm_baseline}. Column~(3) is
identical to Table~\ref{tab:nm_baseline}. $\text{Population}$ is the ZIP code's
own resident population. Quarter fixed effects throughout; standard errors,
in parentheses, are clustered at the county level.
\end{minipage}

\end{table}


\begin{figure}[htbp]
\centering
\includegraphics[width=\textwidth]{fig_nm_zip_noncannabis_placebo_legaleventstudy.png}
\caption{Border coefficient by calendar quarter, non-cannabis sectors, New Mexico ZIP-quarter panel.
$\text{Border}_z$ interacted with each calendar quarter in the sample, anchored to New Mexico's own
legalization date; ZIP-clustered 95\% confidence intervals. Alcohol retail shows fewer quarters because
tax-confidentiality suppression leaves too few ZIP codes reporting to estimate every quarter's interaction.}
\label{fig:noncannabis_eventstudy}
\end{figure}

\begin{figure}[htbp]
\centering
\includegraphics[width=\textwidth]{fig_nm_zip_eventstudy_50_100mi.png}
\caption{Dynamic border premium under the 50-mile and 100-mile border definitions.  Each dot is border-by-quarter coefficients that measures the difference in revenue outcomes of border ZIP codes with interiors. Reference quarter 2022Q2; county-clustered 95\% confidence intervals.}
\label{fig:eventstudy_50_100mi}
\end{figure}


\begin{figure}[htbp]
\centering
\includegraphics[width=\textwidth]{fig_nm_dynamic_legalcontrol.png}
\caption{Dynamic border premium with a legal-border control group, New Mexico. The treated group is border ZIP codes, defined as in the baseline. The control group is New Mexico ZIP codes that satisfy the same rule for the Colorado or Arizona line; ZIP codes near neither kind of line are dropped. Treated ZIP codes: 38; legal-border control ZIP codes: 19.}
\label{fig:dynamic_legalcontrol}
\end{figure}

\begin{figure}[htbp]
\centering
\includegraphics[width=\textwidth]{fig_nm_doseresponse.png}
\caption{Dose-response of dispensary outcomes to distance from the nonlegal-state line, New Mexico. Omitted reference is category 0-25 miles, county-clustered 95\%
confidence intervals.}
\label{fig:nm_doseresponse}
\end{figure}

\begin{figure}[htbp]
\centering
\includegraphics[width=\textwidth]{fig_nm_nonlegalpop_doseresponse.png}
\caption{Dose-response of the New Mexico border premium to nearby nonlegal-state population. Each point is $\beta_k$ from $\log(y_{zt}) = \alpha + \sum_k \beta_k\,\mathbbm{1}[\text{bin}_z=k] + \gamma\log(\text{Population}_z) + \delta_t + \varepsilon_{zt}$, where the bins group New Mexico ZIP codes by the number of Texas, Oklahoma, or Utah residents living within 75 miles of the ZIP centroid; the reference bin is ZIP codes with essentially none. 116-ZIP panel, quarter fixed effects, county-clustered 95\% confidence intervals.}
\label{fig:nm_nonlegalpop_doseresponse}
\end{figure}



\subsection{Cross-State Entry-Regime Comparisons}
\label{sec:crossstate}

\begin{figure}[htbp]
\centering
\includegraphics[width=\textwidth]{fig1_southwest_map.png}
\caption{Dispensary density and the nonlegal-state border, New Mexico and
Colorado, 2025. Shading is dispensaries per 1,000 residents in 2025. Red outlines mark counties that contain at least one
border-exposed dispensary (19 in New Mexico, 17 in Colorado).}
\label{fig:setting}
\label{map}
\end{figure}

\begin{table}[htbp]
\centering
\caption{Cross-state entry-regime comparison: county-level sample, New Mexico and Colorado}
\label{tab:sumstats_state}
\footnotesize
\begin{tabular}{lcc}
\toprule
 & New Mexico & Colorado \\
\midrule
Entry regime & Free entry & Local option / capped \\
First recreational retail sales & 1 April 2022 & 1 January 2014 \\
\addlinespace
Counties (total) & 33 & 64 \\
\quad Border counties (geography) & 11 & 23 \\
\addlinespace
Counties in estimation sample & 32 & 42 \\
County--quarters & 360 & 530 \\
\quad Border county--quarters & 118 & 149 \\
\addlinespace
\multicolumn{3}{l}{\textit{County-quarter means}} \\
\quad Revenue per capita (\$) & 66.5 & 153.1 \\
\quad Dispensaries per 1,000 residents & 0.432 & 0.438 \\
\quad Revenue per dispensary (\$000s) & 213.2 & 524.8 \\
\quad County population & 65,955 & 125,571 \\
\quad Nonlegal pop.\ within 75 mi, border counties (000s) & 194.5 & 37.7 \\
\bottomrule
\end{tabular}
\begin{minipage}{0.86\textwidth}
\vspace{0.4em}\footnotesize
\textit{Notes:} County-quarter panel, 2022Q2 through 2025Q2, the sample used for the cross-state entry-regime comparison. Colorado's own market opened in 2014, but it is restricted to 2022Q2 onward.New Mexico is coded free entry; Colorado is coded locally capped because it delegates licensing to local governments and, as of 2019, 48.4 percent of its counties banned recreational cannabis businesses outright \citep{PayanBrownSong2021}.
\end{minipage}
\end{table}

\begin{table}[htbp]
\centering
\caption{Border $\times$ Restricted interaction: New Mexico vs. Colorado}
\label{tab:regime_interaction}
\footnotesize
\begin{tabular}{lc}
\toprule
Outcome & NM+CO \\
\midrule
Revenue per capita & $-0.458$ \\
 & (0.388) \\
Dispensary density & $-0.932^{*}$ \\
 & (0.488) \\
Revenue per store & $+0.465$ \\
 & (0.314) \\
\bottomrule
\end{tabular}
\begin{minipage}{0.98\textwidth}
\vspace{0.5em}

\fontsize{9pt}{11pt}\selectfont
\textit{Notes:} Each row is the interaction coefficient
$\theta$ from a separate regression for the named outcome; the number in
parentheses directly below it is the standard error, clustered at the county
level. Restricted$_s$ = 1 for Colorado and 0 for New Mexico. Both Border$_c$ and Restricted$_s$ are
time-invariant in this sample. $\theta$ is therefore interpreted as the
difference in the estimated border premium between Colorado and New Mexico,
rather than as an unqualified causal treatment effect.
$^{*}p<0.10$, $^{**}p<0.05$, $^{***}p<0.01$.
\end{minipage}
\end{table}

\begin{table}[htbp]
\centering
\caption{Model-implied dispensary-revenue benchmark, Colorado border counties, 2024}
\label{tab:counterfactual}
\footnotesize
\begin{tabular}{lc}
\toprule
Quantity & Value \\
\midrule
Dispensary-density interaction coefficient $\theta$ & $-0.932$ \\
Free-entry benchmark multiplier $e^{-\theta}$ & 2.54 \\
\midrule
Actual 2024 revenue, CO border counties & \$324.2 M \\
Free-entry-benchmark counterfactual revenue (NM) & \$823.0 M \\
Model-implied additional dispensary revenue & +\$498.8 M \\
95\% county-cluster bootstrap interval & (\$-1.5 M, +\$2{,}526.0 M) \\
\bottomrule
\end{tabular}
\begin{minipage}{0.98\textwidth}
\vspace{0.5em}

\fontsize{9pt}{11pt}\selectfont
\textit{Notes:} The benchmark uses the \emph{dispensary-density} interaction
($\theta=-0.932$, Table~\ref{tab:regime_interaction}), the outcome for which the
difference across entry regimes is statistically significant. It asks how much
more dispensary revenue Colorado's border counties would have generated in 2024
if their density border premium had matched New Mexico's rather than the
restricted-entry premium, holding revenue per store at its observed Colorado
level. The multiplier $e^{-\theta}$ is applied to each Colorado border county's
own 2024 dispensary revenue and summed; New Mexico is the free-entry benchmark
state and contributes only to the estimation of $\theta$. Colorado is restricted
to 2022Q2 onward (Section~\ref{sec:crossstate}). The bootstrap interval resamples
counties with replacement within each of the two states (500 replicates),
refitting the density regression each time; it is wide and reaches down to zero
because $\theta$ is only marginally significant ($p\approx0.06$). The dollar
figure should therefore be read as an illustrative benchmark rather than an
identified estimate of foregone revenue.
\end{minipage}
\end{table}

\section*{Declarations and ethics statements}

\noindent \textbf{Data availability statement.} This study uses only publicly available data. New Mexico cannabis sales and licensed-establishment records are published by the New Mexico Cannabis Control Division; Colorado county-level marijuana sales and licensee records are published by the Colorado Department of Revenue; population, geographic, and business-pattern data are from the U.S. Census Bureau and the SimpleMaps U.S. ZIP Code Database. The code that constructs the analysis panels and reproduces all tables and figures is available from the author on request.

\smallskip
\noindent \textbf{Competing interests.} The author has no competing interests to declare that are relevant to the content of this article.

\smallskip
\noindent \textbf{Funding.} The author received no specific funding for this work.
\newpage
\bibliographystyle{chicago}


\bibliography{bibliography.bib}

\end{document}

